\documentclass[12pt,a4paper]{article}
\usepackage[UTF8]{ctex}
\usepackage{amsmath,amssymb,amsthm}
\usepackage{geometry}
\usepackage{bm}
\usepackage{hyperref}
\usepackage{tikz}
\usetikzlibrary{arrows,arrows.meta,positioning,calc,shapes.geometric}

\geometry{margin=2.5cm}

\title{凸的概念详解：从基础到应用}
\author{第12章抓取与操作补充说明}
\date{}

\newtheorem{theorem}{定理}
\newtheorem{definition}{定义}
\newtheorem{lemma}{引理}
\newtheorem{example}{例子}
\newtheorem{proposition}{命题}

\begin{document}

\maketitle

\tableofcontents
\newpage

\section{概述}

\textbf{凸性}（convexity）是数学、优化和几何中的一个基本概念，在接触运动学、机器人抓取和操作中起着关键作用。理解凸性对于分析可行运动集合、接触约束和抓取稳定性至关重要。

\section{凸集（Convex Set）}

\subsection{定义}

\begin{definition}[凸集]
集合$C \subseteq \mathbb{R}^n$是\textbf{凸集}（convex set），如果对于任意两点$x_1, x_2 \in C$和任意$\lambda \in [0, 1]$，都有：
\begin{equation}
\lambda x_1 + (1-\lambda) x_2 \in C
\label{eq:convex_set}
\end{equation}
\end{definition}

\textbf{几何意义}：连接集合中任意两点的线段完全在集合内。

\subsection{直观理解}

\begin{center}
\begin{tikzpicture}[scale=1.5]
% 凸集
\begin{scope}
\fill[blue!20] (0,0) ellipse (2 and 1.5);
\draw[thick,blue] (0,0) ellipse (2 and 1.5);
\filldraw[red] (-1,0.5) circle (2pt) node[below left] {$x_1$};
\filldraw[red] (1,-0.5) circle (2pt) node[below right] {$x_2$};
\draw[thick,green] (-1,0.5) -- (1,-0.5);
\node[blue] at (0,1.8) {凸集};
\end{scope}

% 非凸集
\begin{scope}[xshift=6cm]
\fill[red!20] (0,0) -- (1.5,1) -- (2,0) -- (1.5,-1) -- (0,0);
\draw[thick,red] (0,0) -- (1.5,1) -- (2,0) -- (1.5,-1) -- cycle;
\filldraw[red] (0.3,0.2) circle (2pt) node[below left] {$x_1$};
\filldraw[red] (1.7,-0.2) circle (2pt) node[below right] {$x_2$};
\draw[thick,green,dashed] (0.3,0.2) -- (1.7,-0.2);
\node[red] at (1,1.8) {非凸集};
\node[red] at (1,-1.3) {线段超出集合};
\end{scope}
\end{tikzpicture}
\end{center}

\subsection{例子}

\textbf{凸集的例子}：
\begin{itemize}
\item 线段、直线、射线
\item 圆、椭圆、球
\item 半空间、多面体
\item 整个空间$\mathbb{R}^n$
\end{itemize}

\textbf{非凸集的例子}：
\begin{itemize}
\item 星形（有"凹陷"）
\item 月牙形
\item 字母"C"的形状
\end{itemize}

\section{凸组合（Convex Combination）}

\subsection{定义}

\begin{definition}[凸组合]
给定$k$个点$x_1, x_2, \ldots, x_k \in \mathbb{R}^n$，它们的\textbf{凸组合}定义为：
\begin{equation}
x = \sum_{i=1}^k \lambda_i x_i
\end{equation}
其中系数$\lambda_i$满足：
\begin{enumerate}
\item \textbf{非负性}：$\lambda_i \geq 0$ 对所有$i = 1, \ldots, k$
\item \textbf{归一性}：$\sum_{i=1}^k \lambda_i = 1$
\end{enumerate}
\end{definition}

\subsection{两个点的凸组合}

对于两个点$x_1$和$x_2$：
\begin{equation}
x = \lambda x_1 + (1-\lambda) x_2, \quad 0 \leq \lambda \leq 1
\end{equation}

\textbf{几何意义}：$x$在连接$x_1$和$x_2$的线段上。

\begin{center}
\begin{tikzpicture}[scale=2]
\filldraw[red] (0,0) circle (2pt) node[below left] {$x_1$};
\filldraw[blue] (3,1) circle (2pt) node[below right] {$x_2$};
\draw[thick,green] (0,0) -- (3,1);
\filldraw[purple] (1.5,0.5) circle (2pt) node[above] {$\lambda=0.5$};
\filldraw[purple] (0.75,0.25) circle (1.5pt) node[below] {$\lambda=0.25$};
\filldraw[purple] (2.25,0.75) circle (1.5pt) node[above] {$\lambda=0.75$};
\node[below] at (1.5,-0.3) {凸组合：连接两点的线段};
\end{tikzpicture}
\end{center}

\subsection{三个点的凸组合}

对于三个点$x_1, x_2, x_3$：
\begin{equation}
x = \lambda_1 x_1 + \lambda_2 x_2 + \lambda_3 x_3, \quad \lambda_i \geq 0, \sum \lambda_i = 1
\end{equation}

\textbf{几何意义}：$x$在由三个点形成的三角形内部（包括边界）。

\begin{center}
\begin{tikzpicture}[scale=2]
\filldraw[red] (0,0) circle (2pt) node[below left] {$x_1$};
\filldraw[blue] (3,0) circle (2pt) node[below right] {$x_2$};
\filldraw[green] (1.5,2) circle (2pt) node[above] {$x_3$};
\fill[purple!20] (0,0) -- (3,0) -- (1.5,2) -- cycle;
\draw[thick,green] (0,0) -- (3,0) -- (1.5,2) -- cycle;
\filldraw[purple] (1.5,0.7) circle (1.5pt);
\node[purple] at (1.5,0.4) {凸组合区域};
\node[below] at (1.5,-0.3) {三角形内部（包括边界）};
\end{tikzpicture}
\end{center}

\section{凸包（Convex Hull）}

\subsection{定义}

\begin{definition}[凸包]
给定集合$S \subseteq \mathbb{R}^n$，$S$的\textbf{凸包}（convex hull）$\text{conv}(S)$是包含$S$的\textbf{最小凸集}，即：
\begin{equation}
\text{conv}(S) = \left\{\sum_{i=1}^k \lambda_i x_i \middle| x_i \in S, \lambda_i \geq 0, \sum_{i=1}^k \lambda_i = 1, k \in \mathbb{N}\right\}
\end{equation}
\end{definition}

\textbf{直观理解}：用橡皮筋套住所有点，橡皮筋收缩后的形状就是凸包。

\subsection{例子}

\begin{center}
\begin{tikzpicture}[scale=1.5]
% 原始点集
\begin{scope}
\filldraw[red] (0,0) circle (2pt);
\filldraw[red] (1,0.5) circle (2pt);
\filldraw[red] (2,0) circle (2pt);
\filldraw[red] (1.5,1.5) circle (2pt);
\filldraw[red] (0.5,1) circle (2pt);
\filldraw[red] (1,2) circle (2pt);
\node[below] at (1,-0.5) {原始点集};
\end{scope}

% 凸包
\begin{scope}[xshift=5cm]
\filldraw[red] (0,0) circle (2pt);
\filldraw[red] (1,0.5) circle (2pt);
\filldraw[red] (2,0) circle (2pt);
\filldraw[red] (1.5,1.5) circle (2pt);
\filldraw[red] (0.5,1) circle (2pt);
\filldraw[red] (1,2) circle (2pt);
\fill[blue!20] (0,0) -- (2,0) -- (1.5,1.5) -- (1,2) -- (0.5,1) -- cycle;
\draw[thick,blue] (0,0) -- (2,0) -- (1.5,1.5) -- (1,2) -- (0.5,1) -- cycle;
\node[blue] at (1,0.8) {凸包};
\node[below] at (1,-0.5) {最小凸集};
\end{scope}
\end{tikzpicture}
\end{center}

\section{凸锥（Convex Cone）}

\subsection{定义}

\begin{definition}[凸锥]
集合$C \subseteq \mathbb{R}^n$是\textbf{凸锥}（convex cone），如果：
\begin{enumerate}
\item 对于任意$x \in C$和$\lambda \geq 0$，有$\lambda x \in C$（锥性）
\item 对于任意$x_1, x_2 \in C$，有$x_1 + x_2 \in C$（凸性）
\end{enumerate}
\end{definition}

\textbf{等价定义}：$C$是凸锥当且仅当对于任意$x_1, x_2 \in C$和$\lambda_1, \lambda_2 \geq 0$，有：
\begin{equation}
\lambda_1 x_1 + \lambda_2 x_2 \in C
\end{equation}

\subsection{几何意义}

\begin{itemize}
\item 凸锥以原点为顶点
\item 如果$x$在锥中，那么从原点到$x$的整条射线都在锥中
\item 锥可以延伸到无穷远
\end{itemize}

\begin{center}
\begin{tikzpicture}[scale=1.5]
% 凸锥
\fill[blue!20] (0,0) -- (2,1) -- (2,2) -- (0,1) -- cycle;
\draw[thick,blue] (0,0) -- (2,1);
\draw[thick,blue] (0,0) -- (0,1);
\draw[thick,blue,dashed] (2,1) -- (2,2);
\draw[thick,blue,dashed] (0,1) -- (2,2);
\filldraw[red] (0,0) circle (2pt) node[below left] {原点};
\node[blue] at (1,0.5) {凸锥};
\node[below] at (1,-0.3) {以原点为顶点};
\end{tikzpicture}
\end{center}

\section{正线性组合（Positive Linear Combination）}

\subsection{定义}

\begin{definition}[正线性组合]
给定$k$个向量$x_1, x_2, \ldots, x_k \in \mathbb{R}^n$，它们的\textbf{正线性组合}（positive linear combination）或\textbf{非负线性组合}定义为：
\begin{equation}
x = \sum_{i=1}^k \lambda_i x_i, \quad \lambda_i \geq 0
\end{equation}
\end{definition}

\textbf{与凸组合的区别}：
\begin{itemize}
\item 凸组合：$\lambda_i \geq 0$且$\sum \lambda_i = 1$（有界）
\item 正线性组合：$\lambda_i \geq 0$（无界，可以任意大）
\end{itemize}

\subsection{正张成（Positive Span）}

\begin{definition}[正张成]
向量集合$\{x_1, \ldots, x_k\}$的\textbf{正张成}（positive span）定义为：
\begin{equation}
\text{pos}(\{x_1, \ldots, x_k\}) = \left\{\sum_{i=1}^k \lambda_i x_i \middle| \lambda_i \geq 0\right\}
\end{equation}
\end{definition}

\textbf{几何意义}：这是以原点为顶点的凸锥。

\section{凸性的性质}

\subsection{基本性质}

\begin{proposition}[凸集的交集]
如果$C_1$和$C_2$是凸集，那么$C_1 \cap C_2$也是凸集。
\end{proposition}

\textbf{推广}：任意多个凸集的交集仍然是凸集。

\begin{proposition}[凸集的并集]
凸集的并集不一定是凸集。
\end{proposition}

\textbf{例子}：两个不相交的圆，它们的并集不是凸集。

\subsection{半空间是凸集}

\begin{proposition}
半空间$H = \{x | a^T x \geq b\}$是凸集。
\end{proposition}

\textbf{证明}：对于$x_1, x_2 \in H$和$\lambda \in [0, 1]$：
\begin{align}
a^T (\lambda x_1 + (1-\lambda) x_2) &= \lambda a^T x_1 + (1-\lambda) a^T x_2 \\
&\geq \lambda b + (1-\lambda) b = b
\end{align}

因此$\lambda x_1 + (1-\lambda) x_2 \in H$。

\subsection{多面体是凸集}

\begin{proposition}
多面体（多个半空间的交集）是凸集。
\end{proposition}

\textbf{证明}：多面体是多个半空间的交集，而半空间是凸集，凸集的交集是凸集。

\section{在接触运动学中的应用}

\subsection{可行Twist集合是凸的}

在接触约束中，可行twist集合：
\begin{equation}
\mathcal{V} = \{V_A | F_i^T (V_A - V_{j(i)}) \geq 0 \text{ 对所有 } i\}
\end{equation}

\textbf{为什么是凸的？}
\begin{itemize}
\item 每个约束$F_i^T (V_A - V_{j(i)}) \geq 0$定义了一个半空间
\item 半空间是凸集
\item 多个半空间的交集是凸集
\item 因此可行twist集合是凸集
\end{itemize}

\subsection{CoR区域是凸的}

在平面问题中，可行CoR区域是凸的。

\textbf{原因}：
\begin{itemize}
\item 每个接触贡献一个约束区域
\item 约束区域的交集是凸的
\item 因此可行CoR区域是凸的
\end{itemize}

\subsection{形状闭合的条件}

形状闭合要求接触wrench的凸包包含原点在其内部。

\textbf{等价条件}：接触wrench正张成整个wrench空间。

\section{凸函数（Convex Function）}

\subsection{定义}

\begin{definition}[凸函数]
函数$f: \mathbb{R}^n \to \mathbb{R}$是\textbf{凸函数}（convex function），如果对于任意$x_1, x_2 \in \mathbb{R}^n$和$\lambda \in [0, 1]$，有：
\begin{equation}
f(\lambda x_1 + (1-\lambda) x_2) \leq \lambda f(x_1) + (1-\lambda) f(x_2)
\label{eq:convex_function}
\end{equation}
\end{definition}

\textbf{几何意义}：函数图像上任意两点之间的线段在函数图像的上方。

\begin{center}
\begin{tikzpicture}[scale=2]
% 凸函数
\draw[thick,blue] (0,2) .. controls (1,1) and (2,0.5) .. (3,1);
\draw[thick,green,dashed] (0,2) -- (3,1);
\filldraw[red] (0,2) circle (1.5pt) node[above left] {$(x_1, f(x_1))$};
\filldraw[red] (3,1) circle (1.5pt) node[below right] {$(x_2, f(x_2))$};
\node[blue] at (1.5,0.2) {凸函数};
\node[green] at (1.5,1.5) {线段在函数上方};
\end{tikzpicture}
\end{center}

\subsection{例子}

\textbf{凸函数}：
\begin{itemize}
\item $f(x) = x^2$
\item $f(x) = |x|$
\item $f(x) = e^x$
\item $f(x) = -\log x$（$x > 0$）
\end{itemize}

\textbf{非凸函数}：
\begin{itemize}
\item $f(x) = -x^2$
\item $f(x) = \sin x$
\end{itemize}

\section{凸优化（Convex Optimization）}

\subsection{凸优化问题}

\textbf{标准形式}：
\begin{align}
\text{最小化} \quad & f(x) \\
\text{使得} \quad & g_i(x) \leq 0, \quad i = 1, \ldots, m \\
& h_j(x) = 0, \quad j = 1, \ldots, p \\
& x \in C
\end{align}

其中：
\begin{itemize}
\item $f$是凸函数
\item $g_i$是凸函数
\item $h_j$是仿射函数（线性函数）
\item $C$是凸集
\end{itemize}

\subsection{凸优化的优势}

\begin{enumerate}
\item \textbf{全局最优}：任何局部最优解都是全局最优解
\item \textbf{高效算法}：存在多项式时间的求解算法
\item \textbf{对偶理论}：强对偶性成立
\end{enumerate}

\section{在机器人学中的应用}

\subsection{接触约束}

\begin{itemize}
\item 可行twist集合是凸的
\item 可行CoR区域是凸的
\item 这简化了分析和优化
\end{itemize}

\subsection{抓取规划}

\begin{itemize}
\item 形状闭合测试可以转化为凸优化问题
\item 抓取质量度量通常是凸函数
\item 可以使用凸优化方法求解
\end{itemize}

\subsection{路径规划}

\begin{itemize}
\item 自由空间（如果凸）的路径规划更简单
\item 可以使用凸优化方法
\end{itemize}

\section{总结}

\subsection{关键概念}

\begin{enumerate}
\item \textbf{凸集}：连接任意两点的线段在集合内

\item \textbf{凸组合}：系数非负且和为1的线性组合

\item \textbf{凸包}：包含集合的最小凸集

\item \textbf{凸锥}：以原点为顶点的凸集，可以延伸到无穷

\item \textbf{正线性组合}：系数非负的线性组合（形成凸锥）

\item \textbf{凸函数}：函数图像上任意两点之间的线段在图像上方
\end{enumerate}

\subsection{重要性质}

\begin{enumerate}
\item 半空间是凸集
\item 多个半空间的交集是凸集（多面体）
\item 凸集的交集是凸集
\item 凸集的并集不一定是凸集
\end{enumerate}

\subsection{在接触运动学中的重要性}

\begin{enumerate}
\item 可行运动集合是凸的
\item 这简化了分析和优化
\item 可以使用凸优化方法
\item 保证了某些优化问题的全局最优解
\end{enumerate}

\end{document}

